% Part: sets-functions-relations
% Chapter: arithmetization
% Section: From N to Z
% Hindi translation (hi-Deva-IN) of Open Logic commit 9620cc73f9c8e0ad003c514a5d3748f29611c4c0.

\documentclass[../../../include/open-logic-section]{subfiles}


\begin{document}

\olfileid[hi]{sfr}{arith}{int}
\olsection{$\Nat$ से $\Int$ तक}

यहाँ दो मूल अवलोकन हैं:
	\begin{enumerate}
		\item प्रत्येक पूर्णांक को $n - m$ के रूप में लिखा जा सकता है, जहाँ $n, m \in\Nat$।
		\item व्यंजक $n - m$ में निहित जानकारी को क्रमित युग्म $\tuple{n, m}$ में भी समान रूप से रखा जा सकता है।
	\end{enumerate}
हम पहले ही जानते हैं कि प्राकृत संख्याओं के क्रमित युग्म $\Nat^2$ के
!!{element} हैं। और हम मान रहे हैं कि~$\Nat$ को समझते हैं। अतः इन दोनों
अवलोकनों पर आधारित एक आरंभिक सुझाव यह है: \emph{पूर्णांकों को प्राकृत
संख्याओं के क्रमित युग्म मानें}।

वास्तव में यह सुझाव अत्यधिक सरल है। हम निश्चय ही चाहते हैं कि
$0- 2 = 4 - 6$ हो। किंतु स्पष्टतः $\tuple{0, 2 }\neq \tuple{4, 6}$। इसलिए
हम केवल यह नहीं कह सकते कि $\Nat^2$ ही पूर्णांकों का समुच्चय है।

पिछली समस्या का व्यापकीकरण करने पर हमें यह चाहिए:
	$$a - b = c - d \text{ तभी और केवल तभी जब }a + d = c + b$$
(यह स्पष्ट होना चाहिए कि पूर्णांकों को इसी प्रकार \emph{व्यवहार करना चाहिए}:
दोनों पक्षों में बस $b$ और $d$ जोड़ दीजिए।) इस व्यवहार को सुनिश्चित करने का
सरल ढंग क्रमित युग्मों के बीच तुल्यता संबंध~$\Intequiv$ इस प्रकार परिभाषित
करना है:
$$\tuple{ a, b } \Intequiv \tuple{c, d}\text{ तभी और केवल तभी जब }a + d = c + b$$
अब हमें दिखाना है कि यह एक तुल्यता संबंध है।
\begin{prop} $\Intequiv$ एक तुल्यता संबंध है।
\end{prop}
	\begin{proof}
	हमें दिखाना है कि $\Intequiv$ स्वतुल्य, सममित और संक्रामक है।
	
	\emph{स्वतुल्यता:} स्पष्टतः $\tuple{a, b} \Intequiv \tuple{a, b}$, क्योंकि $a + b = b + a$।
	
	\emph{सममिति:} मान लीजिए $\tuple{a, b} \Intequiv \tuple{c, d}$, अतः $a + d = c + b$। तब $c + b = a + d$, इसलिए $\tuple{c, d} \Intequiv \tuple{a, b}$।
	
	\emph{संक्रामकता:} मान लीजिए $\tuple{a, b} \Intequiv \tuple{c, d}\Intequiv \tuple{m, n}$। तब $a + d = c + b$ और $c + n = m + d$। अतः $a + d + c + n = c + b + m + d$, और इसलिए $a + n = m + b$। फलतः $\tuple{a, b } \Intequiv \tuple{m, n}$।
\end{proof}

अब इस तुल्यता संबंध से तुल्यता वर्ग बनाए जा सकते हैं:

\begin{defn}
		पूर्णांक, प्राकृत संख्याओं के क्रमित युग्मों के~$\Intequiv$ के अधीन तुल्यता
		वर्ग हैं; अर्थात $\Int = \equivclass{\Nat^2}{\Intequiv}$।
\end{defn}

इस प्रकार नियत की गई परिभाषा पर कई अलग-अलग \emph{दार्शनिक} प्रतिक्रियाएँ
हो सकती हैं। उन प्रतिक्रियाओं पर विचार करने से पहले कुछ तकनीकी विवरणों को
आगे बढ़ाना उपयोगी होगा।

पूर्णांक क्या हैं यह बताने के बाद उन पर मूल फलन और संबंध परिभाषित करने होंगे।
$\equivrep{m, n}{\Intequiv}$ से~$\Intequiv$ के अधीन वह तुल्यता वर्ग सूचित
करें जिसका !!a{element} $\tuple{m, n}$ है।\footnote{ध्यान दें:
\olref[sfr][rel][eqv]{def:equivalenceclass} में दिए गए संकेत का उपयोग करते
हुए इसी वस्तु को $\equivrep{\tuple{m,n}}{\Intequiv}$ लिखते हैं। किंतु उसे पढ़ना
थोड़ा कठिन है।} अर्थात:
	$$\equivrep{m, n}{\Intequiv} = \Setabs{\tuple{a, b} \in \Nat^2}{\tuple{a, b}\Intequiv \tuple{m, n}}$$
अब कुछ परिभाषाएँ दें:
	\begin{align*}
		\equivrep{a, b}{\Intequiv} + \equivrep{c, d}{\Intequiv} &= \equivrep{a + c, b + d}{\Intequiv}\\
		\equivrep{a, b}{\Intequiv} \times \equivrep{c, d}{\Intequiv} &= \equivrep{a c + b  d, a  d + b c}{\Intequiv}\\
		\equivrep{a, b}{\Intequiv} \leq \equivrep{c, d}{\Intequiv} &\text{ तभी और केवल तभी जब }a + d \leq b + c
	\end{align*}	
(प्रचलित रीति के अनुसार सूत्रों को पढ़ना आसान बनाने के लिए मैं `$ab$' को
`$(a \times b)$' के स्थान पर लिख रहा हूँ।) अब यह सुनिश्चित करना होगा कि ये
परिभाषाएँ वैसा ही व्यवहार करें जैसा उन्हें \emph{करना चाहिए}। इसका अर्थ
विस्तार से लिखना और जाँचना श्रमसाध्य है; विवरण \olref[check]{sec} में दिए
गए हैं। संक्षेप में: सब कुछ ठीक काम करता है।

एक अंतिम बात शेष है। हमने पूर्णांकों का निर्माण प्राकृत संख्याओं से किया
है। इसका अर्थ होगा कि प्राकृत संख्याएँ \emph{स्वयं पूर्णांक नहीं हैं}। इसके
दार्शनिक महत्त्व पर हम \olref[ref]{sec} में लौटेंगे। किंतु तकनीकी रूप से
हमें प्राकृत संख्याओं को पूर्णांक \emph{मानकर} बरतने का कोई ढंग चाहिए। विचार
सरल है: प्रत्येक $n \in \Nat$ के लिए बस नियत करें कि
$n_\Int = \equivrep{n, 0}{\Intequiv}$। हमें पुष्टि करनी है कि यह परिभाषा
उचित व्यवहार करती है; अर्थात किन्हीं भी $m, n
\in \Nat$ के लिए
\begin{align*}
	(m + n)_\Int &= m_\Int + n_\Int\\
	(m \times n)_\Int &= m_\Int \times n_\Int\\
	m \leq n &\liff m_\Int \leq n_\Int
\end{align*}
यह सब पर्याप्त सीधा है। उदाहरण के लिए, इनमें दूसरी समानता दिखाने हेतु
प्राकृत संख्याओं के गुणों का उपयोग करके इस प्रकार तर्क कर सकते हैं:
%\begin{proof}
%	Helping myself to the behaviour of the natural numbers, evidently:
	\begin{align*}
%		(m + n)_\Int  = \equivrep{m + n, 0}{\Intequiv} = \equivrep{m + n, 0 + 0}{\Intequiv} = \equivrep{m, 0}{\Intequiv} + \equivrep{n, 0}{\Intequiv} = m_\Int + n_\Int\\
		(m \times n)_\Int  &= \equivrep{m \times n, 0}{\Intequiv} \\
		&= \equivrep{m \times n + 0 \times 0, m \times 0 + 0 \times n}{\Intequiv} \\
		&= \equivrep{m, 0}{\Intequiv} \times \equivrep{n, 0}{\Intequiv} \\
		&= m_\Int \times n_\Int
	\end{align*}
%\end{proof}
अन्य दोनों शर्तों की पुष्टि अभ्यास के रूप में छोड़ते हैं।
\begin{prob}
	दिखाइए कि $(m + n)_\Int = m_\Int + n_\Int$ और
	$m \leq n \liff m_\Int \leq n_\Int$, किन्हीं भी $m, n \in \Nat$ के लिए।
\end{prob}
\end{document}
